\begin{table}
\caption{Formal evaluation matrix for research hypotheses (H1-H4) and research questions (RQ1-RQ6).}
\label{tab:hypotheses_summary}
\begin{tabular}{lllll}
\toprule
Type & ID & Statement & Empirical Findings & Verdict \\
\midrule
Hypothesis & H1 & Adaptive HQNN achieves competitive predictive accuracy and ROC-AUC compared with Fixed HQNN. & Across all 3 medical benchmarks, no statistically significant difference was detected in the paired Wilcoxon analysis (all p >= 0.05). Holdout Test AUC: WDBC 0.9907, Heart 0.9524, Diabetes 0.8144. & SUPPORTED \\
Hypothesis & H2 & Adaptive HQNN significantly reduces average inference-time quantum layer usage compared with Fixed HQNN (L=2.0). & Estimated inference-time quantum-layer reductions observed across all benchmarks: WDBC 98.8%, Heart 87.8%, and Diabetes 79.6% (average inference quantum layers per sample: WDBC 0.02, Heart 0.24, Diabetes 0.41). & SUPPORTED \\
Hypothesis & H3 & Adaptive HQNN provides a measurable predictive-performance versus inference-time quantum-workload trade-off. & Demonstrates a measurable trade-off across benchmarks: reduces estimated inference-time quantum layers by 79.6% to 98.8% with observed CV ROC-AUCs of 0.9923 (WDBC), 0.8138 (Heart), and 0.7982 (Diabetes). & SUPPORTED \\
Hypothesis & H4 & The adaptive quantum strategy generalizes across multiple diverse medical disease modalities. & Validated across oncology (WDBC, 30 features), cardiology (Cleveland Heart, 13 features), and endocrinology (Pima Diabetes, 8 features) using leakage-free inner-CV routing. & SUPPORTED \\
Research Question & RQ1 & How does Adaptive HQNN compare with classical baselines in accuracy and AUC? & Performs competitively with Classical MLP across datasets (CV Accuracy: WDBC 96.92% vs 97.80%; Heart 76.02% vs 83.45%; Diabetes 73.94% vs 77.03%; CV ROC-AUC: WDBC 0.9923 vs 0.9947; Heart 0.8138 vs 0.8967; Diabetes 0.7982 vs 0.8452). & ANSWERED \\
Research Question & RQ2 & How does Adaptive HQNN compare with Fixed HQNN in predictive performance? & Across 5-fold cross-validation, no statistically significant difference was detected in the paired Wilcoxon analysis (all p >= 0.05). Holdout Test Accuracies: WDBC (96.49% vs 97.37%), Heart (85.25% vs 85.25%), Diabetes (74.68% vs 72.08%). & ANSWERED \\
Research Question & RQ3 & By what fraction does adaptive routing reduce inference-time quantum circuit evaluations? & Achieves estimated inference-time quantum-layer reductions of 98.8% on WDBC, 87.8% on Heart Disease, and 79.6% on Diabetes relative to fixed depth-2 execution. & ANSWERED \\
Research Question & RQ4 & Does dynamic layer selection improve quantum efficiency without statistically significant predictive-performance difference? & Yes; paired Wilcoxon signed-rank tests confirm no statistically significant difference was detected in the paired Wilcoxon analysis for ROC-AUC between Adaptive and Fixed HQNN (all p >= 0.05) while reducing estimated inference quantum workload by 79.6% to 98.8%. & ANSWERED \\
Research Question & RQ5 & Which clinical features are associated with different quantum-routing paths? & Surrogate routing analysis shows biomarker distribution variations across paths (e.g., higher dispersion in plasma glucose, BMI, and ST depression on Path 1 and Path 2 cohorts). & ANSWERED \\
Research Question & RQ6 & Can classical MLPs achieve identical efficiency without quantum variational layers? & Adaptive Classical Control (M4) achieves similar efficiency via classical dynamic depth, with observed CV F1-scores of 94.48% (WDBC), 72.90% (Heart), and 65.04% (Diabetes), compared to Adaptive HQNN CV F1-scores of 95.88%, 73.49%, and 66.27%. & ANSWERED \\
\bottomrule
\end{tabular}
\end{table}
